% !TEX root = ../Monografia.tex
\addcontentsline{toc}{chapter}{Resumo}

\begin{center}
\huge{{\bf Resumo}}
\vspace{2cm}
\end{center}

O monitoramento de desempenho em aplicações móveis depende, na prática, de ferramentas externas que exigem configuração especializada e interpretação manual de dados, uma barreira que restringe sua adoção ao longo do desenvolvimento cotidiano. Este trabalho apresenta o \texttt{collector\_flutter}, um pacote Dart/Flutter de código aberto que incorpora coleta de métricas, análise heurística e visualização em tempo real diretamente na aplicação monitorada, sem dependências nativas adicionais para as métricas Dart e com \textit{platform channels} para CPU e bateria em Android.

O protótipo adota arquitetura de três camadas (\textit{Data}, \textit{Domain} e \textit{Presentation}) baseada nos princípios da \textit{Clean Architecture}, coletando métricas de renderização (FPS, \textit{jank}, percentis P50/P95/P99), uso de memória RSS e \textit{heap}, tráfego HTTP, uso de CPU por processo, nível de bateria e eventos customizados do desenvolvedor. Um motor heurístico analisa essas métricas a cada ciclo de coleta e gera recomendações de otimização classificadas por severidade, exibidas em um painel visual integrado. A exportação da sessão em JSON e CSV é disponibilizada via \texttt{ExportService}.

O desenvolvimento seguiu um ciclo de vida iterativo estruturado em quatro fases: revisão bibliográfica, especificação de requisitos, implementação do protótipo e avaliação experimental. A validação qualitativa em quatro cenários controlados (\textit{rebuilds} intensivos, carga de rede, alocação de memória e carga combinada) confirmou a operacionalidade do sistema e a pertinência das heurísticas implementadas. A suíte de 48 testes automatizados cobre os módulos de análise, recomendação, exportação e os \textit{data sources} de CPU e bateria.

\textbf{Palavras-chave}: Flutter. Monitoramento de desempenho. Profiling. Coleta de métricas em tempo real. Eficiência energética. Sustentabilidade digital. Clean Architecture.


\clearpage
\thispagestyle{empty}
\cleardoublepage
